% !TeX program = pdflatex
% =============================================================================
% Lektionsplan (Diskussionsgrundlage, NICHT das Arbeitsblatt selbst)
% Normalkraft -- UZH Praktikum I, Sara Romer Urban, KS im Lee, HS2026
% Klasse: 2e
% =============================================================================
\documentclass[11pt,a4paper]{article}

\usepackage[T1]{fontenc}
\usepackage[utf8]{inputenc}
\usepackage[ngerman]{babel}
\usepackage{lmodern}
\usepackage[a4paper,margin=2.2cm]{geometry}
\usepackage{array,booktabs,tabularx,xltabular}
\usepackage{enumitem}
\usepackage{xcolor}
\usepackage{amsmath}
\usepackage{siunitx}
\sisetup{per-mode=fraction,fraction-command=\frac}
\usepackage{url}
\Urlmuskip=0mu plus 1mu
\def\UrlBreaks{\do\/\do-\do_}

\definecolor{titleblue}{HTML}{183A6B}
\definecolor{accentblue}{HTML}{2E6F9E}
\definecolor{groupteal}{HTML}{1B7F72}
\definecolor{darkgray}{HTML}{4A4A4A}
\definecolor{warnred}{HTML}{9E2E2E}

\setlength{\parindent}{0pt}
\setlength{\parskip}{0.5em}

\newcommand{\Methode}[1]{{\small\color{groupteal}\textbf{#1}}}

\usepackage{titlesec}
\titleformat{\section}{\Large\bfseries\color{titleblue}}{\thesection}{0.6em}{}
\titlespacing*{\section}{0pt}{1.2em}{0.5em}

\begin{document}
\sloppy

\begin{center}
{\Huge\bfseries\color{titleblue}Lektionsplan: Normalkraft}\\[0.4em]
{\large Doppellektion (2$\times$45 Min.) -- Diskussionsgrundlage}
\end{center}
\vspace{0.5em}

\noindent\textbf{Klasse:} 2e (09.09.2026)
\noindent\textbf{Lehrperson:} Loris Keller

\vspace{0.3em}
\noindent\textcolor{darkgray}{Dieses Dokument ist \textbf{keine} Reinschrift des Arbeitsblatts, sondern die
Planungsgrundlage dafür: Lernziele, Aufbau der Lektion und die gewählte
Didaktik-Methode pro Phase (aus \texttt{Methodensammlung\_Physikunterricht.pdf}).
Nach Absprache wird daraus das \texttt{.tex}-Arbeitsblatt/Aufgabenblatt lokal in
VS Code erstellt. Der Aufbau folgt dem von Loris vorgeschlagenen Ablauf:
Federwaage-Repetition $\to$ Tisch/kognitiver Konflikt $\to$ elastische
Unterlage als sichtbarer Beleg $\to$ Definition der Normalkraft $\to$
Klotz-Übung $\to$ schiefe Ebene (grafisch, drei Winkel) $\to$
Hangabtriebskraft-Formeln $\to$ Anwendungsaufgabe.}

% =========================================================================
\section*{1. Abgrenzung und Lernziele}

\textbf{Im Fokus:} Einführung der Normalkraft als Reaktionskraft
(inkl. mikroskopischer Ursache: elektrische Abstossung zwischen Atomen/
Molekülen der Oberfläche, PhysikLibre 4.1.5), danach Anwendung der
bereits bekannten grafischen Kraftzerlegung auf die schiefe Ebene
(PhysikLibre 4.8.1/4.8.2). Hangabtriebskraft und Normalkraft-Komponente
werden \textbf{ausschliesslich grafisch} bestimmt und mit dem
Kraft-Massstab ausgemessen -- Sinus/Cosinus sind der Klasse noch nicht
bekannt.

\textbf{Bereits bekannt (nur Repetition, keine Neueinführung):}
Kräftegleichgewicht als Vektorgleichung, Kraft als vektorielle Grösse,
grafische Kraftzerlegung mit dem Kräfteparallelogramm (alle: letzte
Doppellektion, Kraftaddition \& Kraftzerlegung); die Feder-Masse-Skizze
mit Feder- und Gewichtskraft im Gleichgewicht (Gewichtskraft \&
Federkraft); $F=m\cdot a$ (Vorwissen).

\textbf{Bewusst nicht Teil dieser Einheit:} das dritte Newtonsche
Gesetz/Wechselwirkungsgesetz (Aktio-Reaktio) -- eigene Lektion
\texttt{wechselwirkung/} (2e, 16.09.2026, eine Woche später); die
Normalkraft wird deshalb hier rein phänomenologisch eingeführt (Oberfläche
übt eine sichtbare Gegenkraft senkrecht zu sich aus), ohne den Begriff
Aktio-Reaktio oder das Körper-auf-Unterlage-Gegenstück zu benennen.
Ebenfalls nicht Teil: die rechnerische/trigonometrische Bestimmung von
Hangabtriebskraft und Normalkraft-Komponente (PhysikLibre 4.8.3 --
Sinus/Cosinus den SuS noch nicht bekannt); Reibungskraft im Detail -- bei
der Abschlussaufgabe wird \glqq reibungsfrei\grqq{} explizit
vorausgesetzt. Ebenfalls nicht Teil: Kräfte für gekrümmte Bahnen/
Zentripetalkraft (PhysikLibre 4.9), Kräfte in beschleunigten
Bezugssystemen (PhysikLibre 4.10).

\begin{itemize}[leftmargin=1.3cm,itemsep=0.3em]
  \item[LZ1] die Normalkraft als Reaktionskraft einer Oberfläche
  erklären, die senkrecht zu ihr wirkt und mikroskopisch durch die
  elektrische Abstossung zwischen den Atomen/Molekülen der Oberfläche
  entsteht (PhysikLibre 4.1.5).
  \item[LZ2] in verschiedenen Situationen Gewichtskraft und Normalkraft
  miteinander vergleichen und begründen, wann sie gleich gross sind und
  wann nicht.
  \item[LZ3] die Gewichtskraft auf einer schiefen Ebene grafisch in
  Hangabtriebskraft (parallel zur Ebene) und Normalkraft-Komponente
  (senkrecht zur Ebene) zerlegen, für unterschiedliche Neigungswinkel,
  inkl. korrekt eingezeichneter Normalkraft.
  \item[LZ4] Hangabtriebskraft und Normalkraft-Komponente rein grafisch
  bestimmen -- Hilfslinien parallel/normal zur Ebene, Kräfteparallelogramm
  (hier ein Rechteck) mit der Gewichtskraft als Diagonale, und die beiden
  Seiten mit dem Kraft-Massstab ausmessen (PhysikLibre 4.8.2; ohne
  Sinus/Cosinus).
  \item[LZ5] mithilfe der grafisch bestimmten Hangabtriebskraft und
  $F=m\cdot a$ die Beschleunigung und daraus (mit bekannter Kinematik) die
  Endgeschwindigkeit eines reibungsfrei den Hang hinunterrutschenden
  Körpers berechnen.
\end{itemize}

% =========================================================================
\section*{2. Aufbau der Doppellektion}

\begin{xltabular}{\textwidth}{@{}p{2.6cm}X p{1.2cm}@{}}
\toprule
\textbf{Phase} & \textbf{Inhalt \& Methode} & \textbf{Zeit} \\
\midrule
\endhead

Repetition &
\textbf{Federwaage mit Gewicht.} Nochmals fragen: welche Kräfte wirken
hier, wie sieht das Kräftegleichgewicht aus? Direkter Rückgriff auf
Gewichtskraft/Federkraft und Kraftaddition/-zerlegung. Plenumsgespräch. &
5' \\[0.4em]

Aktivierung &
\textbf{Tisch mit Gegenstand -- kognitiver Konflikt.} Frage: welche
Kräfte wirken hier, und warum bewegt sich der Gegenstand nicht durch den
Tisch? SuS merken: die Gewichtskraft allein erklärt das nicht -- es muss
eine weitere Kraft geben. \Methode{C3 Think-Pair-Share} -- zuerst allein,
dann zu zweit, danach Plenum. &
5' \\[0.4em]

Phänomen &
\textbf{Elastische Unterlage als sichtbarer Beleg.} Bild/Demo einer
verformbaren Unterlage mit aufliegender Masse
(\texttt{Normalkraft\_elastische\_oberfläche.png}): die Verformung macht
die Gegenkraft sichtbar, bevor sie formal benannt wird. Kurzer Impuls,
Plenum. &
5' \\[0.4em]

Theorie &
\textbf{Definition Normalkraft.} Reaktionskraft senkrecht zur Oberfläche;
mikroskopische Ursache: elektrische Abstossung zwischen den Atomen/
Molekülen der Oberfläche (PhysikLibre 4.1.5,
\texttt{atoms-repelling-normal-force.jpg}). Bewusst rein phänomenologisch
(LZ1) -- das Körper-auf-Unterlage-Gegenstück und der Begriff
Aktio-Reaktio bleiben der Wechselwirkungs-Lektion (16.09.2026)
vorbehalten. Lehrervortrag, geleitet. &
6' \\[0.4em]

Übung &
\textbf{Klotz-Aufgabe.} Bild mit Klotz in verschiedenen Situationen
(\texttt{Normalkraft\_klotz.png}): Gewichtskraft und Normalkraft jeweils
vergleichen (gleich gross? grösser? kleiner?) (LZ2). \Methode{C1
Sortieraufgabe}, danach Plenumsabgleich mit Lösungsbild
(\texttt{Normalkraft\_klotz\_loesung.png}). &
10' \\[0.4em]

Übergang &
\textbf{Repetition Kraftzerlegung.} Kurze Erinnerung an das
Kräfteparallelogramm aus der letzten Lektion -- Überleitung zur schiefen
Ebene. Plenum. &
3' \\[0.4em]

Theorie/ Konstruktion &
\textbf{Schiefe Ebene: Hangabtriebskraft grafisch ermitteln.} Vorlage
\texttt{schiefe\_ebene\_1.png} zeigt dieselbe Bergziege auf einer
vereisten Bergwand bei drei Neigungswinkeln $\alpha_1{=}20^\circ$,
$\alpha_2{=}30^\circ$, $\alpha_3{=}60^\circ$ (Gradzahlen vorgeschlagen, im Bild
nicht fix vorgegeben). Am ersten Winkel $\alpha_1$ gemeinsam an der Tafel:
Hilfslinien parallel/normal zur Ebene, Kräfteparallelogramm mit der
Gewichtskraft als Diagonale -- wird zum Rechteck, da die Richtungen
senkrecht aufeinander stehen. Die beiden Rechtecksseiten
($F_{G\parallel}$, $F_{G\perp}$) sowie $F_N$ werden direkt mit dem
Kraft-Massstab ausgemessen, \textbf{keine Formel} nötig (PhysikLibre
4.8.2) (LZ3, LZ4). Lehrervortrag, gemeinsame Tafelkonstruktion. &
10' \\[0.4em]

Übung &
\textbf{Selbständige Zerlegung für $\alpha_2$ und $\alpha_3$.} Auf
derselben Vorlage konstruieren die SuS nun eigenständig für die beiden
verbleibenden Winkel Hangabtriebskraft, Normalkraft-Komponente und
Normalkraft, und lesen alle Kräfte mit dem Massstab ab (LZ3, LZ4).
Feinraster und Massstab wie beim Kraftaddition-Thema (siehe Abschnitt 3).
Kontrolle mit den Lösungsbildern
(\texttt{schiefe\_ebene\_1\_loesung\_kräftezerlegung.png},
\texttt{schiefe\_ebene\_1\_loesung\_normalkraft.png}, die bereits alle
drei Winkel zeigen). \Methode{M9 Produktives Üben} -- durch den Vergleich
der drei Winkel beobachten die SuS selbst, wie sich Hangabtriebskraft und
Normalkraft-Komponente mit der Steigung verändern. Einzel-/Partnerarbeit. &
15' \\[0.4em]

Übung/ Anwendung &
\textbf{Die Bergziege auf dem Eis -- Endgeschwindigkeit.} Dieselbe
Bergziege rutscht reibungsfrei bei $\alpha_2{=}30^\circ$ (Masse
$m=\SI{80}{\kilo\gram}$, Hanglänge $s=\SI{10}{\meter}$, Start aus der
Ruhe) die vereiste Bergwand hinunter. Hangabtriebskraft wie zuvor grafisch
konstruieren und mit dem Massstab ausmessen (erwartungsgemäss
$F_{HA}\approx\SI{390}{\newton}$), daraus mit $F_{HA}=m\cdot a$ die
Beschleunigung ($a\approx\SI{4.9}{\meter\per\second\squared}$) und mit
$v^2=2\cdot a\cdot s$ die Endgeschwindigkeit
($v\approx\SI{9.9}{\meter\per\second}$) bestimmen (LZ5). Einzelarbeit,
Plenumsabgleich. &
16' \\[0.4em]

Abschluss &
\textbf{Zusammenfassung und Ausblick aufs Aufgabenblatt.} Kurze
Zusammenfassung, Hinweis auf das separate Aufgabenblatt für weitere
Übung. Plenum. &
5' \\

\bottomrule
\end{xltabular}

\textcolor{darkgray}{\small Summe: 80 Min. -- 10 Min. Puffer innerhalb der
90 Min. der Doppellektion.}

% =========================================================================
\section*{3. Grafische Aufgaben: Feinraster und Massstab (verbindlich)}

Wie beim Kraftaddition/Kraftzerlegung-Thema erhält jede grafische
Konstruktionsaufgabe -- sowohl in der Lektion als auch im Aufgabenblatt --
ein Feld mit feinem Punkt- oder Liniengitter sowie eine explizit
angegebene Massstabsangabe (z.\,B. \glqq 1\,cm $\hat=$ 1\,N\grqq{}), sodass
die gegebene Gewichtskraft exakt in der angegebenen Länge einzuzeichnen
ist und die konstruierten Komponenten (Hangabtriebskraft,
Normalkraft-Komponente) nachmessbar bleiben.

% =========================================================================
\section*{4. Offene Punkte zur Besprechung}

\begin{itemize}[leftmargin=1.2cm,itemsep=0.4em]
  \item \textbf{Vorgeschlagene Gradzahlen $\alpha_1,\alpha_2,\alpha_3$:}
  Die Vorlage \texttt{schiefe\_ebene\_1.png} zeigt die drei Neigungswinkel
  nur symbolisch, ohne Gradzahlen. Ich habe $20^\circ$/$30^\circ$/$60^\circ$
  vorgeschlagen (steigende Reihenfolge, $30^\circ$ zusätzlich für die
  Bergziege-Endgeschwindigkeitsaufgabe wiederverwendet). Bitte bestätigen
  oder anpassen, falls andere Werte besser zur gezeichneten Steilheit der
  drei Ebenen passen.
  \item \textbf{Zeitbudget Klotz-Übung und schiefe Ebene:} Beide Phasen
  beinhalten reales Zeichnen/Vergleichen/Ausmessen -- falls das
  erfahrungsgemäss länger dauert als veranschlagt, ist die Übung
  \glqq Selbständige Zerlegung für $\alpha_2$ und $\alpha_3$\grqq{} (15')
  der Kandidat zum Kürzen (z.\,B. nur $\alpha_2$, $\alpha_3$ als
  Hausaufgabe/im Aufgabenblatt nachholen), da die Bergziege-Aufgabe auf
  denselben Fertigkeiten aufbaut und im Aufgabenblatt ohnehin weitere
  Übung dazu folgt.
\end{itemize}

% =========================================================================
\section*{5. Anschlussdokument: Aufgabenblatt}

Wie bei den vorherigen Themen entsteht zusätzlich ein separates
Aufgabenblatt mit weiteren Übungsaufgaben zur Normalkraft und zur schiefen
Ebene -- inklusive weiterer grafischer Konstruktionen mit Feinraster und
Massstab (siehe Abschnitt 3) sowie zusätzlicher Rechenaufgaben zu
Hangabtriebskraft und Endgeschwindigkeit.

\end{document}
